% !TeX program = lualatex
% !TeX encoding = utf8
% !TeX spellcheck = uk_UA
% !TeX root = ../main.tex



%=========================================================
\Opensolutionfile{answer}[\currfilebase/\currfilebase-Answers]
\Writetofile{answer}{\protect\section*{\nameref*{\currfilebase}}}
\chapter{Випромінювання електромагнітних хвиль}\label{\currfilebase}\hypertarget{Radiation}{}
\graphicspath{{\currfilebase/Pictures}}
%=========================================================
\epigraph{\Alexander Погляд, який я наважуюся висловити, розглядає випромінювання як вищий різновид коливань
у~силових лініях, що, як відомо, зв'язують між собою частинки, а~також маси речовини.}{Майкл Фарадей,
\enquote{Думки про променеві коливання} (1846)}

Ми розглянули основні властивості електромагнітних хвиль: їхню природу
й~закономірності поширення. Проте залишається ключове питання: як саме виникають ці хвилі?



%% --------------------------------------------------------
\section{Історична довідка: досліди Герца}\hypertarget{hertz-experiments}{}
%% --------------------------------------------------------

\begin{wrapstuff}[type=figure, o, top=3, width=0.5\linewidth]
    \centering
    \includegraphics[width=\linewidth]{Hertz_Dipole}
    \caption{Вібратор Герца.}\label{pic:Hertz_vibrator}
\end{wrapstuff}

У~1888~році\index{Герца досліди} Генріх Герц\footnote{Генріх Рудольф Герц (1857--1894)~--- німецький фізик.}
уперше експериментально отримав електромагнітні хвилі. У~його досліді джерелом
випромінювання слугував так званий \emphx{вібратор Герца}\index{Герца вібратор}~--- прямолінійний провідник із
невеликим повітряним проміжком посередині (рис.~\ref{pic:Hertz_vibrator}). Завдяки цьому
проміжку двом половинам вібратора можна було надати значних зарядів: коли різниця
потенціалів сягала граничного значення, у повітряному зазорі відбувався пробій
(проскакувала іскра), і~заряди через іонізоване повітря перетікали з~однієї половини
вібратора на іншу~--- в~такому \emphx{відкритому контурі} виникали електромагнітні
коливання. Щоб швидкозмінні струми існували лише у~вібраторі й~не замикалися через джерело
живлення, між вібратором і~джерелом вмикали \emphx{дроселі}. Високочастотні коливання
існують, поки іскра замикає проміжок; загасання ж їх відбувається в~основному не шляхом
джоулевих втрат (як у~закритому коливальному контурі), а~шляхом \emphx{випромінювання
електромагнітних хвиль}.

Для виявлення хвиль Герц застосовував другий (\emph{приймальний}) вібратор. Під дією
змінного електричного поля хвилі, що надходить від випромінювача, електрони в~приймальному
вібраторі здійснюють вимушені коливання~--- у~ньому збуджується швидкозмінний струм. Якщо
розміри приймального вібратора такі самі, як і~в~випромінювального, то частоти власних
коливань у~них збігаються, і~вимушені коливання внаслідок \emphx{резонансу} досягають
помітної величини. Ці коливання Герц реєстрував за проскакуванням іскри в~мікроскопічному
зазорі посередині приймального вібратора або за свіченням мініатюрної газорозрядної трубки,
увімкненої між його половинами.

Герц не лише експериментально довів існування електромагнітних хвиль, а~й~уперше почав
вивчати їхні властивості~--- поглинання й~заломлення в різних середовищах, відбивання від
металевих поверхонь тощо. На досліді вдалося також виміряти довжину хвилі й~оцінити
швидкість електромагнітних хвиль, яка виявилася рівною швидкості світла. Так було експериментально підтверджено
\emphx{електромагнітну теорію світла}, яку Максвелл запропонував ще 1865~р. (розділ~\ref{ElMagWaves}).

Електромагнітна хвиля існує \emph{без джерел поля} в~тому сенсі, що після її
випромінювання електромагнітне поле хвилі вже не пов'язане з~джерелом. Саме цим
електромагнітна хвиля відрізняється від статичних електричного й~магнітного полів, які
не існують у~відриві від джерела.



%% --------------------------------------------------------
\section{Елементарна модель випромінювання}\hypertarget{elementary-radiation-model}{}
%% --------------------------------------------------------


%%% --------------------------------------------------------
\subsection*{Якісна картина}
%%% --------------------------------------------------------


Механізм випромінювання найпростіше зрозуміти з~наочного міркування Дж.~Дж.~Томсона.
Розглянемо електричне поле точкового заряду $q$. У~стані спокою це звичайне
електростатичне поле з~радіальними силовими лініями. Нехай у~момент $t=0$ заряд під дією
зовнішньої сили починає рухатися з~прискоренням $a$, а~через невеликий час $\tau$ дія сили
припиняється, після чого заряд рухається рівномірно зі швидкістю $v_q=a\tau$
(рис.~\ref{pic:VelocityGraph}); нехай $v_q\ll c$, а~$\tau$ дуже мале порівняно з~часом
спостереження $T$.

\begin{SCfigure}[1.75][h!]\centering
	\localinput{velocity-graph.tikz}
	\caption{Швидкість заряду: короткий імпульс прискорення тривалістю $\tau$, після якого
	заряд рухається рівномірно\label{pic:VelocityGraph}}
\end{SCfigure}

Оскільки збурення поля поширюється не миттєво, а~зі швидкістю світла $c$, різні
спостерігачі дізнаються про прискорення заряду в~різний момент часу. Нехай пройшов час
$T\gg\tau$ (рис.~\ref{pic:ElementaryModel}):

\begin{itemize}
	\item \emph{Зовнішній спостерігач} (на відстані, більшій за $cT$) іще нічого не знає
	  про рух заряду: сигнал про початок руху до нього не встиг дійти, і~він бачить поле
	  нерухомого заряду в~його \emph{початковому} положенні.
	\item \emph{Внутрішній спостерігач} (на відстані, меншій за $c(T-\tau)$) уже знає, що
	  прискорення закінчилось і~заряд рухається рівномірно,~--- він бачить поле рівномірно
	  рухомого заряду в~його \emph{поточному} положенні (на відстані $v_qT$ від початкового, бо $T\gg\tau$).
	\item У~проміжній \emph{тонкій сферичній оболонці} товщиною $c\tau$~---
	  \emph{області випромінювання}~--- силові лінії поля мусять \enquote{зшивати} ці два
	  різні поля, утворюючи характерний \emph{злам}.
\end{itemize}

\begin{figure}[h!]\centering
	\localinput{elementary-model.tikz}
	\caption{Зламані силові лінії поля короткочасно прискореного
	заряду: межа між \enquote{старим} і~\enquote{новим} полем\label{pic:ElementaryModel}}
\end{figure}

Саме цей злам силових ліній і~є, по суті, електромагнітна хвиля, що
\enquote{відірвалась} від заряду й~розповсюджується самостійно. Оскільки лінії
неперервні (теорема Гаусса), вони не можуть обірватись на межі~--- і~саме необхідність
\enquote{зшити} два різні поля без розриву ліній породжує додаткову, \emph{поперечну}
складову поля, якої немає ні в~чисто електростатичному, ні в~чисто
\enquote{рівномірно-рухомому} полі окремо.


%%% --------------------------------------------------------
\subsection*{Математичне виведення поля випромінювання}
%%% --------------------------------------------------------




Розглянемо детальніше злам силової лінії, що виходила з~заряду під кутом $\theta$ до
напряму його прискорення (рис.~\ref{pic:ElementaryDerivation}). У~момент $T\gg\tau$ злам
лежить у~тонкій сферичній оболонці $c(T-\tau)<r<cT$; оскільки $\tau\ll T$, його відстань
до початкового положення заряду можна вважати рівною $r=cT$. Цей злам виник під час
прискорення ($0<t<\tau$) і~відтоді \enquote{біжить} від заряду зі швидкістю $c$. Розкладемо
поле $\Efield$ у~зламі на \emph{радіальну} (поздовжню) складову $E_\parallel$ і
\emph{поперечну} складову $E_\perp$.

\begin{figure}[h!]\centering
	\localinput{elementary-model-derivation.tikz}
	\caption{Геометрія зламу силової лінії: розкладання поля на поздовжню $E_\parallel$ і
	поперечну $E_\perp$ складові\label{pic:ElementaryDerivation}}
\end{figure}

За час $T$ заряд змістився на $v_qT$, тож поперечне до променя зміщення лінії поля
дорівнює $v_qT\sin\theta$, а~радіальна товщина зламу~--- $c\tau$. Лінія поля в~зламі йде
вздовж гіпотенузи цього трикутника, тому відношення поперечної й~радіальної складових
поля дорівнює відношенню катетів:
% ~~~~~~~~~~~~~~~~~~~~~~~~~~~~~~~~~~~~~~~~~~~~~~~~~~~~~~~~
\begin{equation}\label{eq:RatioEE}
	\frac{E_\perp}{E_\parallel} = \frac{v_qT\sin\theta}{c\tau}.
\end{equation}
% ~~~~~~~~~~~~~~~~~~~~~~~~~~~~~~~~~~~~~~~~~~~~~~~~~~~~~~~~
Поздовжня складова~--- це просто (майже незмінене) кулонівське поле заряду на відстані
$r=cT$ (у~гауссовій системі одиниць):
% ~~~~~~~~~~~~~~~~~~~~~~~~~~~~~~~~~~~~~~~~~~~~~~~~~~~~~~~~
\begin{equation}\label{eq:EparDef}
	E_\parallel = \frac{q}{(cT)^2} \stackrel{r=cT}{=} \frac{q}{(cT)\,r}.
\end{equation}
% ~~~~~~~~~~~~~~~~~~~~~~~~~~~~~~~~~~~~~~~~~~~~~~~~~~~~~~~~
Оскільки прискорення $a=v_q/\tau$, підставляючи в~\eqref{eq:RatioEE} і~\eqref{eq:EparDef}:
% ~~~~~~~~~~~~~~~~~~~~~~~~~~~~~~~~~~~~~~~~~~~~~~~~~~~~~~~~
\begin{equation}\label{eq:EperpDerivation}
	E_\perp(t) = \frac{q}{(cT)\,r}\cdot\frac{v_qT\sin\theta}{c\tau} =
	\frac1{c^2}\frac{q}{r}\, a\,\sin\theta,
\end{equation}
% ~~~~~~~~~~~~~~~~~~~~~~~~~~~~~~~~~~~~~~~~~~~~~~~~~~~~~~~~
де прискорення потрібно брати саме в~той момент, коли воно й~породило спостережуваний
зараз злам, тобто в~момент $t-r/c$:
% ~~~~~~~~~~~~~~~~~~~~~~~~~~~~~~~~~~~~~~~~~~~~~~~~~~~~~~~~
\begin{equation}\label{eq:EperpResult}
	\tcbhighmath{ E_\perp(t) = \frac1{c^2}\,\frac{q}{r}\,
	a\!\left(t-\frac{r}{c}\right)\sin\theta. }
\end{equation}
% ~~~~~~~~~~~~~~~~~~~~~~~~~~~~~~~~~~~~~~~~~~~~~~~~~~~~~~~~

Формула~\eqref{eq:EperpResult} містить дві ключові риси \emph{будь-якого} випромінювання,
які ми надалі зустрічатимемо знову й~знову:
\begin{enumerate}
	\item Напруженість поля \emph{хвилі} $E_\perp$ спадає як $1/r$, на відміну від
	  електростатичного поля $E_\parallel\propto1/r^2$. Це узгоджується із законом
	  збереження енергії: у~вакуумі хвиля не поглинається, її енергія розподіляється по поверхні сфери ($\propto r^2$),
	  а~густина енергії $\propto E^2$, тому $E^2\cdot r^2 = \const \Rightarrow
	  E\propto1/r$.
	\item $E_\perp$ у~момент часу $t$ визначається прискоренням заряду $a$ не в~цей самий
	  момент, а~в~\emph{більш ранній} момент $t-r/c$~--- \emphx{час запізнення},
	  потрібний хвилі, щоб дістатися точки спостереження на відстані~$r$.
\end{enumerate}



%% --------------------------------------------------------
\subsection*{Випромінювання осцилюючого заряду}
%% --------------------------------------------------------



Нехай тепер заряд $q$ не просто короткочасно прискорюється, а~весь час рухається вздовж
прямої зі змінним прискоренням $a(t)$ поблизу початку координат~--- наприклад, здійснює
гармонічні коливання. Тоді він випромінює електромагнітні хвилі \emph{неперервно},
і~напруженість поля випромінювання в~точці на відстані $r$ від початку координат, як
і~раніше, визначається формулою~\eqref{eq:EperpResult}, але тепер прискорення потрібно
брати в~більш ранній момент $t-r/c$:
% ~~~~~~~~~~~~~~~~~~~~~~~~~~~~~~~~~~~~~~~~~~~~~~~~~~~~~~~~
\begin{equation}\label{eq:OscillatingChargeE}
	E_\perp(t) = \frac{1}{c^2}\,\frac{q}{r}\,a\!\left(t-\frac{r}{c}\right)\sin\theta.
\end{equation}
% ~~~~~~~~~~~~~~~~~~~~~~~~~~~~~~~~~~~~~~~~~~~~~~~~~~~~~~~~

Нехай, зокрема, заряд здійснює гармонічні коливання з~амплітудою $A$ і~частотою $\omega$:
% ~~~~~~~~~~~~~~~~~~~~~~~~~~~~~~~~~~~~~~~~~~~~~~~~~~~~~~~~
\begin{equation*}
	x(t) = A\cos\omega t, \qquad a(t) = \ddot x(t) = -\omega^2 A\cos\omega t.
\end{equation*}
% ~~~~~~~~~~~~~~~~~~~~~~~~~~~~~~~~~~~~~~~~~~~~~~~~~~~~~~~~
Підставляючи це прискорення у~\eqref{eq:OscillatingChargeE}, отримуємо
% ~~~~~~~~~~~~~~~~~~~~~~~~~~~~~~~~~~~~~~~~~~~~~~~~~~~~~~~~
\begin{equation}\label{eq:OscillatingChargeField}
	E_\perp(t) = -\frac{q\omega^2 A}{c^2
	r}\,\cos\omega\!\left(t-\frac{r}{c}\right)\sin\theta.
\end{equation}
% ~~~~~~~~~~~~~~~~~~~~~~~~~~~~~~~~~~~~~~~~~~~~~~~~~~~~~~~~
Отже, поле в~кожній точці простору гармонічно коливається з~тією самою частотою $\omega$:
\emph{осцилюючий заряд випромінює монохроматичну хвилю}.
Формула~\eqref{eq:OscillatingChargeField} справедлива на відстанях $r$, великих порівняно
з~амплітудою коливань $A$ (і, звісно, за умови $v\ll c$).

\begin{Attention}[unbreakable]
	\emph{Монопольного випромінювання не існує!} Повний заряд $Q$ системи стала величина,
	а~поле сферично-симетричного розподілу зовні, за теоремою Гаусса, дорівнює $Q/r^2$
	незалежно від руху всередині, тому \enquote{дихальні} коливання заряду хвилі не породжують.
	Але тоді як випромінює \emph{один} заряд, що коливається?
	Коли заряд коливається біля певного \enquote{центра}, він фактично утворює
	\emphx{диполь} із цим центром~--- і~саме зміна дипольного моменту системи
	визначає випромінювання. Тому найпростішою системою, здатною ефективно випромінювати, є
	\emphx{електричний диполь} зі змінним у~часі дипольним моментом. До нього ми
	й~переходимо.
\end{Attention}



%% --------------------------------------------------------
\section{Дипольне випромінювання}\hypertarget{dipole-radiation}{}
%% --------------------------------------------------------



%% --------------------------------------------------------
\subsection*{Диполь як елементарний випромінювач}
%% --------------------------------------------------------



Розглянемо\index{Дипольне випромінювання} електронейтральну систему з~двох зарядів $+q$ і~$-q$, з~яких додатний~---
нескінченно важкий (і~тому нерухомий), а~від'ємний вільно рухається
(рис.~\ref{pic:DipoleTwoCharges}). Дипольний момент такої системи та його друга похідна за
часом:
% ~~~~~~~~~~~~~~~~~~~~~~~~~~~~~~~~~~~~~~~~~~~~~~~~~~~~~~~~
\begin{equation*}
	\vect p = -q\vect r', \qquad \ddot{\vect p} = -q\dot{\vect v},
\end{equation*}
% ~~~~~~~~~~~~~~~~~~~~~~~~~~~~~~~~~~~~~~~~~~~~~~~~~~~~~~~~
де $\vect v=\dot{\vect r}'$~--- швидкість заряду, а~$\dot{\vect v}$~--- його прискорення.

\begin{figure}[h!]\centering
	\localinput{dipole-two-charges.tikz}
	\caption{Найпростіша модель диполя: важкий нерухомий заряд $+q$ і~легкий рухомий
	$-q$\label{pic:DipoleTwoCharges}}
\end{figure}



%% --------------------------------------------------------
\subsection*{Про запізнюючі потенціали (коротко)}
%% --------------------------------------------------------



Строге виведення полів диполя спирається на розв'язання рівнянь Максвелла через
електромагнітні потенціали $\phi$, $\vect A$ у~калібруванні Лоренца, що приводить до
неоднорідного хвильового рівняння Даламбера з~розв'язками у~вигляді \emphx{запізнюючих
потенціалів}:
% ~~~~~~~~~~~~~~~~~~~~~~~~~~~~~~~~~~~~~~~~~~~~~~~~~~~~~~~~
\begin{equation*}
	\vect A(\vect r, t)= \frac1c \iiint\limits_{V'} \frac{\vect j\!\left(\vect r', t -
	\dfrac{|\vect r-\vect r'|}{c}\right) dV'}{|\vect r-\vect r'|}, \qquad \phi(\vect r, t)
	= \iiint\limits_{V'} \frac{\rho\!\left(\vect r', t - \dfrac{|\vect r-\vect
	r'|}{c}\right) dV'}{|\vect r-\vect r'|}.
\end{equation*}
% ~~~~~~~~~~~~~~~~~~~~~~~~~~~~~~~~~~~~~~~~~~~~~~~~~~~~~~~~

Ці потенціали називають \emphx{запізнюючими}\index{Потенціали запізнюючі}, бо поле в~точці $\vect r$ у~момент $t$
визначається не миттєвим розподілом зарядів і~струмів, а їхнім станом у~\emph{попередній}
момент часу $t-|\vect r-\vect r'|/c$~--- з~урахуванням скінченної швидкості поширення
поля. Формально рівняння Даламбера допускають і~другий розв'язок, зі знаком \enquote{$+$}
в~аргументі (\emphx{випереджувальні потенціали}),~--- але він не має ясного фізичного
сенсу (поле випереджало б причину, яка його породжує) і~відкидається на основі
\emphx{причинності}. Розкладаючи ці інтеграли в~ряд для випадку $r\gg r'$ (система
набагато менша за відстань до спостерігача) і~нерелятивістських швидкостей зарядів ($v\ll
c$), можна довести, що головний член розкладу визначається виключно \emphx{дипольним
моментом} системи $\vect p(t-r/c)$~--- це й~строго обґрунтовує, чому диполь є найпростішим
(і~найважливішим) елементарним випромінювачем; наступні члени розкладу відповідають
магнітно-дипольному, електрично-квадрупольному й~вищим типам випромінювання, зазвичай
значно слабшим за дипольне.



%% --------------------------------------------------------
\subsection*{Поля диполя: ближня й~хвильова зони}
%% --------------------------------------------------------



Нехай\index{Ближня і хвильова зони} дипольний момент змінюється гармонічно, $\vect p = \vect p_0\cos\omega t$. Точне
вираження полів (у~сферичних координатах, з~віссю $z$ уздовж $\vect p_0$) містить доданки
різного порядку за $1/r$, і~його зручно аналізувати окремо у~двох граничних випадках.



\subsubsection*{Ближня зона}
Відстань від осцилятора до точки спостереження мала порівняно 
з~довжиною хвилі:
% ~~~~~~~~~~~~~~~~~~~~~~~~~~~~~~~~~~~~~~~~~~~~~~~~~~~~~~~~
\begin{equation}\label{eq:NearZoneCondition}
	r \ll \frac\lambda{2\pi}.
\end{equation}
% ~~~~~~~~~~~~~~~~~~~~~~~~~~~~~~~~~~~~~~~~~~~~~~~~~~~~~~~~
Тут електричне поле в~кожен момент часу $t$ практично збігається з~полем
\emph{статичного} диполя, момент якого дорівнює миттєвому значенню $p(t)$:
% ~~~~~~~~~~~~~~~~~~~~~~~~~~~~~~~~~~~~~~~~~~~~~~~~~~~~~~~~
\begin{equation}\label{eq:NearZoneE}
	E_r = \frac{2\cos\theta\, p(t)}{r^3}, \qquad E_\theta = \frac{\sin\theta\, p(t)}{r^3},
\end{equation}
% ~~~~~~~~~~~~~~~~~~~~~~~~~~~~~~~~~~~~~~~~~~~~~~~~~~~~~~~~
а~магнітне поле, оскільки $\dot p(t) = Il$ (де $I$~--- еквівалентний струм, $l$~---
довжина еквівалентного елемента струму), збігається з квазістатичним полем цього елемента
струму за законом Біо--Савара--Лапласа:
% ~~~~~~~~~~~~~~~~~~~~~~~~~~~~~~~~~~~~~~~~~~~~~~~~~~~~~~~~
\begin{equation}\label{eq:NearZoneB}
	\Bfield = \frac1{cr^3}\left[\dot{\vect p}(t)\times\vect r\right].
\end{equation}
% ~~~~~~~~~~~~~~~~~~~~~~~~~~~~~~~~~~~~~~~~~~~~~~~~~~~~~~~~
У~ближній зоні поля \emph{не запізнюються} (запізнення $r/c \ll T = 2\pi/\omega$~---
мізерне порівняно з~періодом коливань) і~практично збігаються з~полями відповідного
статичного диполя й~елемента струму~--- випромінювання тут ще \enquote{не встигло
сформуватися}.



\subsubsection*{Хвильова зона}
Відстань, набагато більша за довжину хвилі:
% ~~~~~~~~~~~~~~~~~~~~~~~~~~~~~~~~~~~~~~~~~~~~~~~~~~~~~~~~
\begin{equation}\label{eq:WaveZoneCondition}
	r \gg \frac\lambda{2\pi}.
\end{equation}
% ~~~~~~~~~~~~~~~~~~~~~~~~~~~~~~~~~~~~~~~~~~~~~~~~~~~~~~~~
Тут (рис.~\ref{pic:WaveZone}) залишаються лише члени, що спадають як $1/r$ (решта~---
швидше, і~на великих відстанях нехтовно малі):
% ~~~~~~~~~~~~~~~~~~~~~~~~~~~~~~~~~~~~~~~~~~~~~~~~~~~~~~~~
\begin{equation}\label{eq:WaveZoneFields}
	\tcbhighmath{ E_\theta = B_\phi = \frac{\sin\theta}{c^2r}\,\ddot{p}\!\left(t-\frac{r}c\right), \qquad E_r=E_\phi=B_r=B_\theta=0. }
\end{equation}
% ~~~~~~~~~~~~~~~~~~~~~~~~~~~~~~~~~~~~~~~~~~~~~~~~~~~~~~~~

\begin{SCfigure}[1][h!]\centering
	\includegraphics[width=0.45\linewidth]{wave_zone}
	\caption{Поля $E_\theta$ і~$B_\phi$ у~хвильовій зоні диполя\label{pic:WaveZone}}
\end{SCfigure}

У~хвильовій зоні поля $\Efield$ і~$\Bfield$ поперечні, чисельно рівні одне одному
й~спадають як $1/r$~--- точнісінько ті самі риси плоскої електромагнітної хвилі
(розділ~\ref{ElMagWaves}). Крім того, поле залежить від полярного кута $\theta$: на осі
диполя ($\theta=0,\pi$) поле дорівнює нулю, максимуму ж воно сягає в~екваторіальній
площині ($\theta=\pi/2$)~--- \emph{диполь не випромінює вздовж власної осі}. У~кожній
точці хвильової зони вектори $\Efield$, $\Bfield$ і~$\vect r$ взаємно перпендикулярні
й~утворюють правовинтову трійку.

Формула~\eqref{eq:WaveZoneFields}~--- \emph{загальна}: вона описує випромінювання
будь-якого диполя з~довільним законом зміни моменту $p(t)$. Зокрема, для гармонічного
осцилятора $p(t)=p_0\cos\omega t$ вона дає
% ~~~~~~~~~~~~~~~~~~~~~~~~~~~~~~~~~~~~~~~~~~~~~~~~~~~~~~~~
\begin{equation*}
	E_\theta =
	-\frac{p_0\omega^2}{c^2r}\,\cos\omega\!\left(t-\frac{r}{c}\right)\sin\theta,
\end{equation*}
% ~~~~~~~~~~~~~~~~~~~~~~~~~~~~~~~~~~~~~~~~~~~~~~~~~~~~~~~~
що з~точністю до позначень ($p_0=qA$) збігається
з~формулою~\eqref{eq:OscillatingChargeField}, отриманою вище для осцилюючого заряду. Отже,
\emph{осцилюючий заряд є частковим випадком диполя}~--- і~загальна теорія дипольного
випромінювання природно узагальнює результати попереднього параграфа (рис.~\ref{pic:DipoleRadiationSnapshot}).

\begin{SCfigure}[1][h!]
	\centering
	\includegraphics[width=0.5\linewidth]{dipole_radiation}
	\caption{Електричне поле осцилюючого диполя: у~ближній зоні поле ніби
	\enquote{причеплене} до диполя, у хвильовій~---
	випромінюється}
	\label{pic:DipoleRadiationSnapshot}
\end{SCfigure}


%% --------------------------------------------------------
\subsection*{Діаграма спрямованості}
%% --------------------------------------------------------


Кутову залежність поля випромінювання диполя ($\propto\sin\theta$) зручно
зображати \emphx{полярною діаграмою спрямованості}\index{Діаграма спрямованості} (рис.~\ref{pic:PolarPattern}): довжина
відрізка від початку координат до кривої в~напрямку $\theta$ дорівнює величині поля у~цьому
напрямку (на рисунку~--- два кола). Для інтенсивності ($\propto\sin^2\theta$) крива вужча, але
має той самий вигляд \enquote{бублика}: нулі вздовж осі диполя, максимум в~екваторіальній площині.

\begin{SCfigure}[2][h!]
	\centering
		\localinput{polar-pattern.tikz}
	\caption{Полярна діаграма спрямованості дипольного випромінювача: характерна форма
	\enquote{бублика}\label{pic:PolarPattern}}
\end{SCfigure}



%% --------------------------------------------------------
\subsection*{Поляризація випромінювання}
%% --------------------------------------------------------



З~формули~\eqref{eq:WaveZoneFields} видно, що у~хвильовій зоні відмінна від нуля лише
\emph{одна} компонента електричного поля~--- $E_\theta$: вектор $\Efield$ у~кожній точці
простору лежить у~площині, що проходить через вісь диполя (напрямок $\vect p$) і~напрямок
спостереження $\vect r$. Оскільки ця площина для заданого напрямку спостереження не
змінюється з~часом (диполь коливається вздовж фіксованої осі), випромінена хвиля
виявляється \emphx{лінійно поляризованою}: вектор $\Efield$ у~кожній точці простору
коливається вздовж однієї й~тієї самої прямої, а~не обертається й~не описує еліпс. Це не
просто математична деталь: саме завдяки лінійній поляризації дипольного випромінювання
приймальні антени ефективно вловлюють сигнал лише тоді, коли їх зорієнтовано
\emph{вздовж} площини коливань поля хвилі (див.~п.~\enquote{Опір випромінювання}), а~поляризаційні
світлофільтри здатні пропускати чи затримувати світло певної поляризації (детальніше про
поляризацію~--- розділ~\ref{ElMagWaves}).



%% --------------------------------------------------------
\section{Потужність випромінювання осцилятора}\hypertarget{oscillator-radiated-power}{}
%% --------------------------------------------------------

Густина потоку електромагнітної енергії характеризується вектором Пойнтінга $\vect\Pi$
(розділ~\ref{MaxwellEqns}); потік енергії крізь поверхню сфери радіуса $r$, що оточує
осцилятор, і~називають \emphx{потужністю випромінювання}\index{Потужність випромінювання}. Елемент площі сфери в~сферичних
координатах $d\sigma = r^2\sin\theta\, d\theta\, d\phi$, тому, підставляючи поля
хвильової зони~\eqref{eq:WaveZoneFields} у~$\Pi = \dfrac{c}{4\pi}EB = \dfrac{c}{4\pi}E^2$:
% ~~~~~~~~~~~~~~~~~~~~~~~~~~~~~~~~~~~~~~~~~~~~~~~~~~~~~~~~
\begin{multline}\label{eq:PowerInstant}
	\frac{dW}{dt} = \oiint\limits_\sigma \Pi\, d\sigma =
	\frac{c}{4\pi}\int\limits_0^{2\pi}\!\!\int\limits_0^\pi
	\left[\frac{\sin\theta}{c^2r}\, \ddot{p}\!\left(t-\frac{r}{c}\right)\right]^2
	r^2\sin\theta\, d\theta\, d\phi \\ {}= \frac{\ddot{p}^2\!\left(t-\dfrac{r}{c}\right)}{4\pi c^3}\cdot2\pi\int\limits_0^\pi\sin^3\theta\,
	d\theta.
\end{multline}
% ~~~~~~~~~~~~~~~~~~~~~~~~~~~~~~~~~~~~~~~~~~~~~~~~~~~~~~~~
Оскільки $\displaystyle\int_0^\pi\sin^3\theta\,d\theta = \frac43$, отримуємо
% ~~~~~~~~~~~~~~~~~~~~~~~~~~~~~~~~~~~~~~~~~~~~~~~~~~~~~~~~
\begin{equation}\label{eq:PowerInstantResult}
	\frac{dW}{dt} = \frac{2\ddot{p}^2\!\left(t-\dfrac{r}{c}\right)}{3c^3}.
\end{equation}
% ~~~~~~~~~~~~~~~~~~~~~~~~~~~~~~~~~~~~~~~~~~~~~~~~~~~~~~~~
Тут $\ddot{p}$~--- друга похідна дипольного моменту за часом (тобто \enquote{прискорення}
дипольного моменту). Для гармонічного осцилятора $p=p_0\cos\omega t$, $\ddot{p} =
-\omega^2p_0\cos\omega t$, і, усереднюючи за період ($\langle\cos^2\rangle=1/2$):
% ~~~~~~~~~~~~~~~~~~~~~~~~~~~~~~~~~~~~~~~~~~~~~~~~~~~~~~~~
\begin{equation}\label{eq:PowerAveraged}
	\tcbhighmath{ \left\langle\frac{dW}{dt}\right\rangle = \frac{\omega^4p_0^2}{3c^3}. }
\end{equation}
% ~~~~~~~~~~~~~~~~~~~~~~~~~~~~~~~~~~~~~~~~~~~~~~~~~~~~~~~~

\begin{wrapstuff}[type=figure, o, top=2, width=0.4\linewidth]
    \centering
    \includegraphics[width=\linewidth]{HalfWaveDipole}
	\caption{Реальна півхвильова дипольна антена\label{pic:HalfWaveDipole}}
\end{wrapstuff}
Середня потужність випромінювання осцилятора пропорційна квадрату амплітуди дипольного
моменту $p_0$ і~--- значно різкіше~--- \emph{четвертому степені частоти} $\omega$ (за незмінної амплітуди $p_0$; для антени зі сталою
амплітудою струму $p_0=I_0l/\omega$, і~потужність $\propto\omega^2$). Тому
ефективно випромінюють лише досить швидкі коливання: дроти зі струмом промислової частоти $50$~Гц (довжина хвилі
$6000$~км) практично нічого не випромінюють, а~для радіозв'язку використовують частоти від сотень кілогерц і~вище (довжини
хвиль від кілометрів до сантиметрів), коли розміри антени можна зробити порівнянними з~довжиною хвилі. Реальні
радіоантени часто мають розмір порядку половини довжини хвилі~--- саме такі
\emph{півхвильові диполі} (рис.~\ref{pic:HalfWaveDipole}) найефективніше випромінюють
і~приймають електромагнітні хвилі.



%% --------------------------------------------------------
\subsection*{Розсіювання світла: чому небо блакитне}
%% --------------------------------------------------------



Тим\index{Розсіювання світла (чому небо блакитне)} самим характером залежності випромінювання осцилятора від частоти пояснюється, зокрема,
і~блакитний колір неба (рис.~\ref{pic:BlueSky}). Сонячне світло, що пронизує атмосферу,
розсіюється молекулами повітря, які можна вважати елементарними осциляторами: під впливом
світлової хвилі вони здійснюють \emphx{вимушені} коливання. Оскільки період власних
коливань цих молекул істотно відрізняється від періоду видимого світла (резонанс
відсутній), амплітуда вимушених коливань слабко залежить від частоти хвилі, тому
інтенсивність розсіяного світла, тобто інтенсивність випромінювання цих вимушених
осциляторів, за формулою~\eqref{eq:PowerAveraged}, обернено пропорційна $\lambda^4$
(позаяк $\omega\propto1/\lambda$). Тому короткохвильове (синє) світло розсіюється значно
сильніше, ніж довгохвильове (червоне),~--- і~небо здається блакитним. Фіолетове світло
розсіюється ще сильніше, але воно слабше представлене в~спектрі Сонця, частково поглинається
у~верхніх шарах атмосфери, а~око до нього менш чутливе, тому в~сумі переважає блакитний колір.

\begin{wrapstuff}[type=figure, o, top=2, width=0.5\linewidth]
    \centering
    \includegraphics[width=\linewidth]{BlueSky}
	\caption{Розсіювання сонячного світла молекулами атмосфери~--- причина блакитного
	кольору неба\label{pic:BlueSky}}
\end{wrapstuff}
Та сама залежність $\propto1/\lambda^4$ пояснює і, на перший погляд, протилежне явище~---
\emph{червоне Сонце на заході}. Коли Сонце стоїть низько над горизонтом, його промені,
перш ніж дістатися ока спостерігача, проходять крізь \emph{набагато товщий} шар
атмосфери, ніж опівдні. Синя й~фіолетова складові світла, розсіюючись найсильніше,
встигають майже повністю \enquote{вибути} з~прямого пучка вздовж цього довгого шляху
(розсіюючись убік~--- це і~є та сама блакить, яку ми бачимо, дивлячись не на Сонце,
а~в~інший бік неба); до ока ж у~прямому напрямку доходить переважно те, що лишилося
в~пучку,~--- довгохвильове, червоно-жовтогаряче світло. Отже, блакитне небо вдень
і~червоне Сонце на заході~--- два прояви однієї й~тієї самої фізики \emphx{релеївського
розсіювання}: розсіяне вбік світло~--- блакитне, послаблений розсіюванням прямий пучок~---
червоний.



%% --------------------------------------------------------
\subsection*{Опір випромінювання}
%% --------------------------------------------------------



З~погляду генератора, що живить антену змінним струмом, випромінена антеною енергія~--- це
енергія, яка \emph{безповоротно} йде з~кола, точнісінько так само, як і~енергія, що
розсіюється на звичайному омічному опорі. Тому в~теорії кіл зручно формально ввести
\emphx{опір випромінювання}\index{Опір випромінювання} $R_\text{вип}$~--- такий еквівалентний опір, що виділяв би
в~собі ту саму середню потужність $\langle dW/dt\rangle$ при тому самому амплітудному
струмі $I_0$ в~антені, який реально в~ній тече:
% ~~~~~~~~~~~~~~~~~~~~~~~~~~~~~~~~~~~~~~~~~~~~~~~~~~~~~~~~
\begin{equation*}
	\left\langle\frac{dW}{dt}\right\rangle \equiv \frac12 I_0^2 R_\text{вип}.
\end{equation*}
% ~~~~~~~~~~~~~~~~~~~~~~~~~~~~~~~~~~~~~~~~~~~~~~~~~~~~~~~~
Це дозволяє формально врахувати випромінювання нарівні зі звичайними тепловими втратами
при розрахунку електричних кіл антен, хоча фізична природа цих \enquote{втрат} зовсім
інша: енергія не переходить у~тепло всередині самої антени, а~назавжди йде з~нею у~вигляді
електромагнітної хвилі. Той самий півхвильовий диполь чудово працює й~у~зворотному
режимі~--- як приймальна антена. Характеристики
антени~--- діаграма спрямованості й~поляризаційні властивості~--- однакові й~для
передавання, і~для приймання. Зокрема, приймальна антена найкраще ловить хвилю, коли її вісь
паралельна $\Efield$, і~зовсім не ловить, коли вісь перпендикулярна до $\Efield$.



%% --------------------------------------------------------
\section{Випромінювання рухомого заряду}\hypertarget{moving-charge-radiation}{}
%% --------------------------------------------------------

Скориставшись моделлю \enquote{важкий $+q$~--- легкий рухомий $-q$} (див.
рис.~\ref{pic:DipoleTwoCharges}) і~формулою~\eqref{eq:PowerInstantResult}, підставимо
$\ddot{\vect p} = -q\dot{\vect v}$:
% ~~~~~~~~~~~~~~~~~~~~~~~~~~~~~~~~~~~~~~~~~~~~~~~~~~~~~~~~
\begin{equation}\label{eq:LarmorFormula}
	\tcbhighmath{ \frac{dW}{dt} = \frac{2q^2}{3c^3}\dot{\vect v}^2. }
\end{equation}
% ~~~~~~~~~~~~~~~~~~~~~~~~~~~~~~~~~~~~~~~~~~~~~~~~~~~~~~~~
Тут $\dot{\vect v}$~--- прискорення заряду (друга похідна його радіус-вектора за часом).
Формулу~\eqref{eq:LarmorFormula} називають \emphx{формулою Лармора}\index{Лармора формула}.

\begin{Attention}
	Формула~\eqref{eq:LarmorFormula}~--- один із найважливіших результатів класичної
	електродинаміки: \emph{будь-який} заряд, що рухається з~прискоренням, неминуче
	випромінює електромагнітну енергію, причому (при $v\ll c$) втрачена потужність залежить \emph{лише}
	від заряду й~миттєвого прискорення~--- і~зовсім не залежить від того, чому саме заряд
	прискорюється (електричне поле, магнітне поле, механічний поштовх). Заряд, що
	рухається \emph{рівномірно} й~прямолінійно ($\dot{\vect v}=0$), не випромінює
	взагалі~--- випромінює лише \emph{прискорений} заряд. Сама формула нерелятивістська: при
	$v\sim c$ потужність залежить і~від швидкості (див.\ синхротронне випромінювання нижче).
\end{Attention}



%% --------------------------------------------------------
\subsection*{Гальмівне випромінювання}
%% --------------------------------------------------------



Різке гальмування швидких заряджених частинок~--- поширений і~технологічно важливий прояв
формули Лармора, який називають \emphx{гальмівним випромінюванням}\index{Гальмівне випромінювання} (нім.
\textit{Bremsstrahlung}). Саме на ньому ґрунтується робота рентгенівської трубки: пучок
електронів, розігнаний високою напругою (десятки кіловольт), влучає в~металевий анод
і~різко (за час прольоту відстані порядку розміру атома) гальмується електричним полем атомних
ядер речовини анода. За формулою~\eqref{eq:LarmorFormula} таке величезне (від'ємне)
прискорення неминуче супроводжується випромінюванням; оскільки різні електрони пучка
гальмуються по-різному~--- одні зупиняються практично одразу в~одному зіткненні, інші~---
поступово, за кілька послідовних зіткнень, віддаючи енергію частинами,~--- випромінювання
виходить не на одній частоті, а~суцільним спектром. Саме так рентгенівські трубки й
отримують потрібне для медичної діагностики, дефектоскопії металів та рентгеноструктурного
аналізу випромінювання.



%% --------------------------------------------------------
\subsection*{Синхротронне випромінювання}
%% --------------------------------------------------------



Той самий принцип~--- прискорено рухомий по коловій траєкторії заряд неминуче
випромінює~--- працює не лише для кулонівського притягання, а~й для руху заряджених
частинок по колу в~магнітному полі (де прискорення забезпечує сила Лоренца, а~не
кулонівська сила). На цьому явищі, яке називають \emphx{синхротронним випромінюванням}\index{Синхротронне випромінювання},
ґрунтується робота потужних лабораторних джерел рентгенівського й~ультрафіолетового
випромінювання (синхротронів) для дослідження речовини й~матеріалів; водночас воно ж~---
головне джерело небажаних втрат енергії в~кільцевих прискорювачах заряджених частинок.
Оскільки втрати різко зростають зі зменшенням маси частинки, для легких,
ультрарелятивістських електронів синхротронні втрати на коловій орбіті стають
катастрофічними (для ультрарелятивістської частинки потужність у~$\gamma^4$ разів більша
за ларморівську, $\gamma=W/mc^2$, тому при тій самій енергії вона $\propto m^{-4}$; задача~\ref{prb:SynchrotronRatio}): саме вони обмежили енергію найбільшого кільцевого електрон-позитронного колайдера LEP (близько
$100$~ГеВ на пучок), тому для вищих енергій електронів проєктують \emph{лінійні} прискорювачі, тоді як важчі протони
у~Великому адронному колайдері прискорюють по кільцю до $6{,}8$~ТеВ (проєктне значення $7$~ТеВ).



%% --------------------------------------------------------
\subsection*{Стабільність атома: класична суперечність}
%% --------------------------------------------------------



Формула\index{Стабільність атома} Лармора має й~драматичний історичний наслідок: електрон, що обертається по коловій
орбіті навколо ядра (як у~класичній \enquote{планетарній} моделі атома Резерфорда),
рухається з~доцентровим прискоренням, а~отже, за~\eqref{eq:LarmorFormula}, \emph{неминуче
випромінює} й~безперервно втрачає енергію. Оцінка (задача~\ref{prb:AtomCollapse}) показує,
що класичний електрон мав би \emph{впасти на ядро} за час порядку $10^{-11}$~с~--- атоми
в~такій моделі просто не могли б бути стабільними!

\begin{Historical}
	Суперечність між класичною електродинамікою й~спостережуваною стабільністю атомів
	стала одним з~головних поштовхів до створення квантової механіки: М.~Бор 1913~р. був
	змушений постулювати, що на деяких (\enquote{стаціонарних}) орбітах електрон
	\emph{всупереч} класичній електродинаміці не випромінює~--- цей постулат, який не мав
	жодного пояснення в~межах класичної фізики, згодом отримав строге обґрунтування вже
	в~повній квантовій теорії.
\end{Historical}



%% --------------------------------------------------------
\subsection*{Реакція випромінювання і~парадокси самодії}
%% --------------------------------------------------------



Формула~\eqref{eq:LarmorFormula} описує втрату енергії, але не пояснює, \emph{як саме}
заряд \enquote{відчуває} цю втрату. Якщо заряд випромінює, то на нього має діяти додаткова
сила з~боку його власного поля~--- \emphx{сила реакції випромінювання}\index{Реакції випромінювання сила} (самодія).
У~класичній електродинаміці її записують у~вигляді формули Абрагама--Лоренца:
% ~~~~~~~~~~~~~~~~~~~~~~~~~~~~~~~~~~~~~~~~~~~~~~~~~~~~~~~~
\begin{equation}\label{eq:RadiationReaction}
	\vect F_{\rm rad} = \frac{2q^2}{3c^3}\,\ddot{\vect v}.
\end{equation}
% ~~~~~~~~~~~~~~~~~~~~~~~~~~~~~~~~~~~~~~~~~~~~~~~~~~~~~~~~
Тут $\ddot{\vect v}$~--- \emph{третя} похідна радіус-вектора заряду за часом (похідна
прискорення, його \enquote{ривок}). Через це рівняння руху заряду
з~урахуванням~\eqref{eq:RadiationReaction} стає диференціальним рівнянням \emph{третього}
порядку~--- і~саме звідси виникають класичні парадокси.

\begin{Attention}
	Рівняння з~третьою похідною має, крім фізичних, два неприйнятні типи розв'язків:
	\begin{itemize}
		\item \emphx{самочинний розгін}\index{Самодії парадокси} (runaway): заряд без зовнішніх сил експоненційно
		  збільшує прискорення, тобто енергія нібито береться \enquote{нізвідки};
		\item \emphx{передприскорення} (pre-acceleration): заряд починає реагувати на силу
		  ще до того, як вона ввімкнулася, що в наївному вигляді порушує причинність.
	\end{itemize}
	До цього додається проблема точкового заряду: його власна енергія розбігається,
	і~коректне віднімання нескінченностей вимагає перенормування. У~межах класичної
	електродинаміки ці труднощі так і~не були остаточно подолані~--- вони стали одним із
	кроків до квантової електродинаміки, де реакція випромінювання описується
	випромінюванням реальних фотонів та обміном віртуальними.
\end{Attention}



%% --------------------------------------------------------
\section*{Підсумки}
\phantomsection\addcontentsline{toc}{section}{Підсумки}
%% --------------------------------------------------------

Розділ про випромінювання підбиває важливий підсумок: \emph{будь-який} прискорений заряд неминуче породжує хвилю, що назавжди йде
від нього в~навколишній простір,~--- а~електромагнітне поле остаточно постає перед нами не
як допоміжна математична конструкція для опису сил між зарядами, а~як цілком самостійна,
фізично реальна субстанція, здатна існувати, поширюватись і~переносити енергію \emph{без}
будь-яких зарядів поблизу. Саме на цьому явищі~--- випромінюванні прискорених зарядів ---
ґрунтується вся радіотехніка, вся астрофізика випромінювання зір і~туманностей, і,
зрештою, саме воно~--- а~не якийсь окремий, спеціально влаштований \enquote{носій
світла}~--- відповідає за існування й~поширення видимого світла в~природі.



%% --------------------------------------------------------
\section*{Задачі}



\phantomsection \addcontentsline{toc}{section}{Задачі}
%% --------------------------------------------------------

% задачі цього розділу короткі: не розриваємо їх між сторінками

\begin{problem}\label{prb:DipolePowerFromFlux}
	Електромагнітна хвиля, що випромінюється диполем, поширюється у~вакуумі так, що у~хвильовій зоні на промені,
	перпендикулярному до осі диполя, на відстані $r$ від нього середнє значення густини потоку енергії дорівнює
	$\langle\Pi\rangle$. Знайдіть середню потужність випромінювання диполя.
	\begin{solution}
		при $\theta=\pi/2$ $\langle\Pi\rangle=\dfrac{\langle\ddot p^2\rangle}{4\pi c^3r^2}$, тому $\left\langle\dfrac{dW}{dt}\right\rangle=\dfrac{2\langle\ddot p^2\rangle}{3c^3}=\dfrac{8\pi r^2}{3}\langle\Pi\rangle$.
	\end{solution}
\end{problem}

\begin{problem}\label{prb:RotatingDipole}
	Сталий за модулем електричний дипольний момент $\vect p$ обертають з~кутовою швидкістю $\omega$ навколо осі,
	перпендикулярної до осі диполя, що проходить через його середину. Знайдіть потужність випромінювання диполя.
	\begin{solution}
		$p_x=p\cos\omega t$, $p_y=p\sin\omega t$, $\ddot{\vect p}^{\,2}=p^2\omega^4=\const$, тому $\dfrac{dW}{dt}=\dfrac{2p^2\omega^4}{3c^3}$~--- удвічі більше, ніж у~лінійного осцилятора з~тією самою амплітудою, і~стала в~часі: обертовий диполь~--- це два лінійні осцилятори, зсунуті за фазою на $90^\circ$, пульсації потужності яких взаємно компенсуються.
	\end{solution}
\end{problem}

\begin{problem}\label{prb:CoulombOrbitPower}
	Знайдіть потужність випромінювання нерелятивістської частинки із зарядом $e$ і~масою $m$, що рухається коловою орбітою
	радіуса $R$ у~полі нерухомого точкового заряду $q$.
	\begin{solution}
		$\dot v=\dfrac{v^2}{R}=\dfrac{eq}{mR^2}$, тому $\dfrac{dW}{dt}=\dfrac{2e^2}{3c^3}\left(\dfrac{eq}{mR^2}\right)^2=\dfrac{2e^4q^2}{3c^3m^2R^4}$.
	\end{solution}
\end{problem}

\begin{problem}\label{prb:AtomCollapse}
	\emph{Падіння класичного електрона на ядро.} Електрон рухається навколо протона по коловій орбіті радіуса
	$a_0=0{,}53\cdot10^{-8}$~см і~через випромінювання поступово втрачає енергію. Вважаючи, що в~кожен момент орбіта колова,
	оцініть, за який час електрон упав би на ядро.
	\begin{solution}
		енергія $W=-\dfrac{e^2}{2r}$, втрати $\dfrac{dW}{dt}=-\dfrac{2e^6}{3c^3m^2r^4}$, звідки $\dfrac{dr}{dt}=-\dfrac{4e^4}{3m^2c^3r^2}$ і~$t=\dfrac{m^2c^3a_0^3}{4e^4}\approx1{,}6\cdot10^{-11}$~с.
	\end{solution}
\end{problem}

\begin{problem}\label{prb:BlueRed}
	\emph{Синє й~червоне.} У~скільки разів синє світло ($\lambda=450$~нм) розсіюється молекулами повітря сильніше, ніж червоне
	($\lambda=700$~нм)?
	\begin{solution}
		інтенсивність розсіяння $\propto\lambda^{-4}$, тому $(700/450)^4\approx6$.
	\end{solution}
\end{problem}

\begin{problem}\label{prb:ShortAntenna}
	\emph{Опір випромінювання короткої антени.} Антена~--- прямий провідник довжиною $l=1$~м, по якому тече (майже однаковий
	уздовж провідника) струм з~амплітудою $I_0=1$~А і~частотою $10$~МГц. Знайдіть опір випромінювання й~потужність випромінювання.
	\emph{Вказівка:} амплітуда дипольного моменту $p_0=I_0l/\omega$ (бо $\dot p=Il$); $1$~с/см $=9\cdot10^{11}$~Ом.
	\begin{solution}
		$\langle dW/dt\rangle=\dfrac{\omega^2I_0^2l^2}{3c^3}=\dfrac12I_0^2R_\text{вип}$, тож $R_\text{вип}=\dfrac{2\omega^2l^2}{3c^3}=\dfrac{8\pi^2}{3c}\left(\dfrac l\lambda\right)^2$, у~СІ $R_\text{вип}=80\pi^2(l/\lambda)^2$~Ом $\approx0{,}9$~Ом при $\lambda=30$~м; потужність $\approx0{,}44$~Вт. Коротка антена ($l\ll\lambda$) випромінює погано.
	\end{solution}
\end{problem}

\begin{problem}\label{prb:SynchrotronRatio}
	\emph{Синхротронні втрати електрона й~протона.} Для ультрарелятивістської частинки на коловій орбіті потужність
	випромінювання в~$\gamma^4$ разів більша за ларморівську, де $\gamma=W/mc^2$. У~скільки разів потужність випромінювання
	електрона більша, ніж протона з~тією самою енергією на орбіті того самого радіуса?
	\begin{solution}
		прискорення $v^2/R\approx c^2/R$ однакові, тож $P\propto\gamma^4\propto m^{-4}$, і~$P_e/P_p=(m_p/m_e)^4\approx1{,}1\cdot10^{13}$.
	\end{solution}
\end{problem}

\begin{Questions}
\begin{quizlist}
	\item У~чому полягав дослід Герца 1888~р.? Як була влаштована установка
	  з~випромінювальним та приймальним вібраторами, і~чому в~приймачі виникає резонансний
	  відгук?
	\item Чим електромагнітна хвиля принципово відрізняється від статичних електричного
	  й~магнітного полів у~сенсі зв'язку з~джерелом?
	\item Опишіть якісну картину випромінювання короткочасно прискореного заряду: чим
	  відрізняється поле, яке бачить зовнішній спостерігач, від поля, яке бачить
	  внутрішній, і~що відбувається в~проміжній сферичній оболонці товщиною $c\tau$?
	\item Чому саме необхідність \enquote{зшити} без розриву силові лінії поля до й~після
	  прискорення заряду породжує додаткову поперечну складову поля?
	\item Виведіть формулу поперечної складової поля випромінювання $E_\perp =
	  \dfrac{q}{c^2r}\,a(t-r/c)\sin\theta$, використовуючи подібність трикутників у~зламі
	  силової лінії.
	\item Чому напруженість поля хвилі спадає як $1/r$, на відміну від електростатичного
	  поля $\propto1/r^2$? Як це узгоджується із законом збереження енергії?
	\item Що таке час запізнення і~чому поле випромінювання в~момент $t$ визначається
	  прискоренням заряду не в~цей самий момент, а~в~момент $t-r/c$?
	\item Чому монопольного випромінювання не існує? Як ізольований коливний заряд, попри
	  закон збереження заряду, все ж утворює випромінюючу систему~--- диполь?
	\item Що таке запізнюючі потенціали і~чим вони відрізняються від випереджувальних?
	  Чому саме запізнюючий розв'язок відповідає фізичній причинності?
	\item Чим поле диполя в~ближній зоні ($r\ll\lambda/2\pi$) відрізняється від поля
	  в~хвильовій зоні ($r\gg\lambda/2\pi$)? Чому в~ближній зоні поле практично збігається
	  зі статичним полем диполя того самого моментального значення $p(t)$?
	\item Запишіть поля хвильової зони диполя $E_\theta=B_\phi=\sin\theta\,\ddot{p}(t-r/c)/(c^2r)$. Які риси плоскої електромагнітної хвилі вони відтворюють? Чому
	  диполь не випромінює вздовж власної осі?
	\item Що таке діаграма спрямованості дипольного випромінювача? Чому вона має
	  характерну форму \enquote{бублика}, з~максимумом в~екваторіальній площині?
	\item Чому випромінювання диполя виявляється лінійно поляризованим? Яке практичне
	  значення це має для орієнтації приймальних антен?
	\item Виведіть формулу миттєвої потужності дипольного випромінювання $dW/dt = 2\ddot{p}^2(t-r/c)/(3c^3)$, інтегруючи вектор Пойнтінга по сфері у~хвильовій зоні.
	\item Виведіть середню потужність випромінювання гармонічного осцилятора $\langle
	  dW/dt\rangle = \omega^4p_0^2/(3c^3)$. Чому саме різка залежність від четвертого
	  степеня частоти пояснює перевагу коротких хвиль для ефективного радіовипромінювання?
	\item Поясніть, чому небо блакитне, спираючись на залежність інтенсивності розсіяного
	  світла $\propto1/\lambda^4$. Як та сама залежність пояснює червоний колір Сонця на
	  заході?
	\item Що таке опір випромінювання антени? Чим фізична природа цих \enquote{втрат}
	  відрізняється від звичайних омічних втрат, попри формальну схожість у~розрахунку
	  кіл?
	\item Виведіть формулу Лармора $dW/dt = 2q^2\dot{\vect v}^2/(3c^3)$ для потужності
	  випромінювання довільного прискореного заряду. Чому рівномірно рухомий заряд не
	  випромінює взагалі?
	\item Що таке гальмівне випромінювання (Bremsstrahlung)? Чому воно виходить не на
	  одній частоті, а~суцільним спектром, і~як це використовується в~рентгенівських
	  трубках?
	\item Що таке синхротронне випромінювання? Чому синхротронні втрати особливо
	  катастрофічні для легких ультрарелятивістських частинок, і~чому це змушує будувати
	  найбільші прискорювачі електронів лінійними, а~не кільцевими?
	\item У~чому полягає класична суперечність між формулою Лармора та стабільністю атома?
	  Який історичний вихід із цієї суперечності запропонував Бор?
	\item Що таке сила реакції випромінювання (формула Абрагама--Лоренца)? Які парадокси
	  (самочинний розгін, передприскорення) виникають через рівняння руху третього порядку?
\end{quizlist}
\end{Questions}

\Closesolutionfile{answer}
